% V2第26章思考题26.1--26.13独立概念解答，未独立专家审校。
\section*{第26章\quad 波粒二象性：思考题解析}
\paragraph{思考题 26.1}
霓虹灯在低压气体放电下，电场激发气体原子/
离子，电子跃迁发出离散谱线，是\emph{放电发光}，
一般不能把它近似为对应气体温度的黑体辐射。
熔炉里的高温铁水主要由热平衡近似的
\emph{热辐射}发光，谱形受温度与发射率影响。
\[
\boxed{\text{霓虹灯：放电跃迁谱；热铁水：热辐射为主。}}
\]

\paragraph{思考题 26.2}
人体 \(T\approx310\,\mathrm K\) 的黑体谱峰
由Wien位移定律
\[
\boxed{\lambda_{\max}\simeq
\frac{2.898\times10^{-3}\,\mathrm{m\,K}}{310\,\mathrm K}
\simeq9.35\,\mathrm{\mu m}}
\]
属于热红外，而人眼敏感波段约
\(0.4\)--\(0.7\,\mathrm{\mu m}\)。
可见频段的人体热辐射在此温度下指数抑制，
所以普通人眼在黑暗中看不见；
热成像仪可用红外探测器观测人体热辐射。

\paragraph{思考题 26.3}
房屋开窗像一个多次漫反射的
\emph{空腔入口}。外界射入的可见光经室内
墙壁和家具多次吸收、散射，
回到小窗口并朝观察者方向逸出的光相对较少。
白天外墙被阳光强烈照亮，窗口区域
的亮度与外墙形成高对比度，
看上去近黑。
\[
\boxed{\text{多次吸收和亮度反差使窗口呈黑洞样外观。}}
\]
它不意味着室内无光或窗口是真正温度黑体。

\paragraph{思考题 26.4}
在高温\emph{相同温度}下，
Kirchhoff热辐射定律给
\(\epsilon_\lambda=\alpha_\lambda\)。
烟煤表面通常在可见区吸收率、
发射率较高；洁净白瓷表面通常较低，
故与同温瓷砖相比
\[
\boxed{\text{烟煤涂覆区域通常更亮、红热更明显}}.
\]
严格比较须给出两表面在发光波段的
谱发射率并确认两处温度相同；
若材料烧蚀、温差或背景反光显著则可能不同。

\paragraph{思考题 26.5}
深色花瓣通常较强吸收可见太阳光，
吸收能流导致花瓣温度上升；
过热会使水分蒸发增大，
部分色素与组织更容易受光热损伤。
\[
P_{\rm abs}\simeq\alpha I_{\odot}A,\qquad
P_{\rm rad}\simeq\epsilon\sigma A(T^4-T_{\rm amb}^4).
\]
因此受控日照/室内栽培可能帮助深色花朵保持
外观。**但原题所说具体公园的养护做法，
不能仅由热辐射公式确证**；
还取决于品种、光周期、
花青素和湿度等园艺因素，需实地资料核查。

\paragraph{思考题 26.6}
若替换 \(h'=10^{34}h\)，
对给定频率 \(\nu\)，光量子能量
\(E'=h'\nu\) 巨大；
同一有限几何尺度中的粒子能级间距
\(E_n\propto h^2\) 会大幅增大。
在固定 \(T\) 的黑体模型，
频率谱峰值
\(\nu_m=x_\nu k_BT/h'\) 因而下移
\(10^{34}\) 倍；
总辐射系数
\(\sigma\propto h^{-3}\) 极小
（若同时假定 \(k_B,c\) 不变）。
德布罗意波长 \(\lambda=h/p\) 也会大幅增加。
\[
\boxed{\text{宏观尺度的量子效应可变得显著，
经典日常物理将不再适用。}}
\]
这种只改一个基本常数的反事实世界
会改变原子结构、物质稳定性与温度尺度，
不能把现实世界其他条件无矛盾地全部固定。

\paragraph{思考题 26.7}
入射频率 \(\nu>\nu_0\) 满足光电阈值时，
\[
\boxed{K_{\max}=h\nu-W}.
\]
(1) 固定 \(\nu\) 增加光强，
在量子效率不变、未达饱和时光子通量
增加，\emph{光电子数/光电流增加}，
但 \(K_{\max}\) 不变。
(2) 固定\emph{能流光强} \(I\)
提高 \(\nu\)，单个光子能量增加，
\(K_{\max}\) 增加；但光子数流量
\(\Phi_\gamma=I/(h\nu)\) \emph{减少}。
故若量子效率保持不变，光电子计数通常减少，
而不是必然增加。
量子效率若强烈依赖频率，
实际光电子数也可能非单调。

\paragraph{思考题 26.8}
同一材料不同电子在金属内的初始状态、
能带能量及逃逸路径不同，
在光吸收后到达表面可能经散射损失不同能量。
因此光电子具有一段动能分布，
只有最有利的初始状态且无额外能量损失
才达到
\[
\boxed{K_{\max}=h\nu-W}.
\]
如果表面不洁净，局部功函数也会变化，
使能量分布更宽。
不能把“单色光子能量完全相同”
误解成所有逸出电子具有相同动能。

\paragraph{思考题 26.9}
康普顿波长变化
\[
\boxed{\Delta\lambda=
\frac{h}{m_ec}(1-\cos\theta)}
\]
最大为
\(2h/(m_ec)\simeq4.85\,\mathrm{pm}\)。
对 \(\lambda\sim500\,\mathrm{nm}
=5\times10^5\,\mathrm{pm}\) 的可见光，
\(\Delta\lambda/\lambda\lesssim10^{-5}\)，
极难同束光的原有谱宽及介质弹性散射分开。
在真实物质里电子还可能处于束缚态，
Rayleigh/Thomson等弹性过程占重要地位。
康普顿效应\emph{不是在可见光绝对不存在}，
而是其相对频移很小，难以通常手段分辨。

\paragraph{思考题 26.10}
普通光电效应中，光子能量用于
克服电子逸出功与形成电子动能；
电子动量的反冲可以由
\emph{晶格或原子核整体}承担，
故简化公式只显式保留电子能量关系。
相反，康普顿模型的碰撞对象是
初始静止的近自由单电子，
\emph{没有可忽略反冲的外部晶格}。
因此必须联立
\[
\boxed{E_\gamma+m_ec^2
=E'_\gamma+E'_e},\qquad
\boxed{\boldsymbol p_\gamma
=\boldsymbol p'_\gamma+\boldsymbol p'_e}.
\]
若不保留动量守恒，就无法推导与散射角有关的
康普顿波长位移。

\paragraph{思考题 26.11}
若两种粒子的非相对论动能相同 \(K\)，
动量 \(p=\sqrt{2mK}\)，
德布罗意波长
\[
\boxed{\lambda=\frac{h}{\sqrt{2mK}}}.
\]
质子质量约为电子1836倍，
因此
\[
\boxed{\frac{\lambda_e}{\lambda_p}
=\sqrt{\frac{m_p}{m_e}}\simeq42.8},
\]
即\emph{电子波长更长}。
在高能相对论情况下应使用
\(p c=\sqrt{K(K+2mc^2)}\)。

\paragraph{思考题 26.12}
对固定经典动量 \(p\)，
\(h\to0\) 使
\(\lambda_{\rm dB}=h/p\to0\)，
量子干涉相位迅速振荡，适当粗粒化与
退相干条件下趋向\emph{经典轨迹/几何光学极限}。
不确定关系的下界
\(\Delta x\Delta p\ge\hbar/2\) 亦形式上趋零，
但\emph{不是}单独令 \(h=0\) 就必然获得
所有经典动力学，需同时给定物理尺度。
若另令 \(c\to\infty\)，
\(\gamma=(1-v^2/c^2)^{-1/2}\to1\)，
Lorentz变换在有限速度下趋于Galilei变换，
相对论时间膨胀与长度收缩消失。
\[
\boxed{h\to0:\text{经典作用量极限};
\quad c\to\infty:\text{非相对论时空极限}}.
\]
二者是不同的独立极限。

\paragraph{思考题 26.13}
束缚在分子势阱中的粒子不可能同时
\(\Delta x=0,\Delta p=0\)。
谐振子即使 \(T=0\)，基态仍有
\[
\boxed{E_0=\frac12\hbar\omega>0},
\qquad
\boxed{\langle K\rangle=\langle V\rangle
=\frac14\hbar\omega},
\]
即\emph{量子零点运动}。
但一颗完全自由且非局域的理想粒子
可以取 \(p=0\) 的动量本征态；
其位置波函数在无穷空间不平方可积，
不是局域的“静止小球”。
所以零温并不自动等于局域束缚自由度
机械意义上的完全不运动。
